% =============================================================
% บทที่ 5 - บทสรุปและข้อเสนอแนะ
% =============================================================

% -----------------------------------------------------------
\section{สรุปผลการดำเนินโครงงาน}
% -----------------------------------------------------------

โครงงาน ``ระบบจำแนกและวิเคราะห์ข่าวสารอัตโนมัติด้วยเทคนิคการประมวลผลภาษาธรรมชาติแบบเฉพาะประเภท'' ได้ดำเนินการพัฒนาเสร็จสิ้นตามขอบเขตและข้อกำหนดที่วางไว้ โดยสามารถรวบรวมข่าวสารจากหลากหลายแหล่ง สรุปเนื้อหาและวิเคราะห์อารมณ์ของข่าวด้วย Large Language Model (KKU Gen.AI) จำแนกหมวดหมู่ข่าวภาษาไทยด้วยเทคนิค Thai NLP และแสดงผลแบบเรียลไทม์ผ่าน WebSocket พร้อมระบบแนะนำข่าวสาร ผลการดำเนินงานสรุปเปรียบเทียบกับวัตถุประสงค์ได้ดังตารางที่ \ref{tab:obj_result}

\begin{table}[H]
\centering
\caption{สรุปผลการดำเนินงานเทียบกับวัตถุประสงค์}
\label{tab:obj_result}
\resizebox{\textwidth}{!}{%
\begin{tabular}{|c|p{4.5cm}|p{6cm}|c|}
\hline
\textbf{ข้อ} & \textbf{วัตถุประสงค์} & \textbf{ผลการดำเนินงานที่ได้} & \textbf{สถานะ} \\
\hline
1 & เพื่อศึกษาและวิเคราะห์เทคนิคการดึงข้อมูลข่าวด้วย Web Scraping, RSS Feed และเทคนิค Thai NLP & ศึกษาและประยุกต์ใช้ Playwright สำหรับ JavaScript-rendered sites, ใช้ BeautifulSoup สำหรับ HTML parsing และ PyThaiNLP สำหรับการจำแนกหมวดหมู่ & บรรลุ \\ \hline
2 & เพื่อออกแบบและพัฒนาระบบรวบรวม สรุป และวิเคราะห์ข่าวสารแบบ End-to-End Pipeline & พัฒนาสถาปัตยกรรมแบบ Layered Architecture พร้อมระบบแคชสองระดับ, ระบบ Real-time WebSocket และหน้าเว็บ Single Page Application ได้อย่างสมบูรณ์ & บรรลุ \\ \hline
3 & เพื่อทดสอบและประเมินประสิทธิภาพของระบบในด้านความถูกต้องและความรวดเร็ว & ระบบผ่านการทดสอบ UAT ครบทุกข้อ และตอบสนองคำขอเรียกดูข่าวผ่านแคชได้ในเวลาเฉลี่ยต่ำกว่า 150 มิลลิวินาที & บรรลุ \\
\hline
\end{tabular}%
}
\end{table}

จากตารางที่ \ref{tab:obj_result} การดำเนินโครงงานบรรลุวัตถุประสงค์ที่ตั้งไว้ครบทั้ง 3 ข้อ

% -----------------------------------------------------------
\section{ข้อจำกัดของระบบ}
% -----------------------------------------------------------

\subsection{ข้อจำกัดด้านแหล่งข้อมูลภายนอก}

ข้อจำกัดประการแรกสืบเนื่องมาจากโครงสร้างหน้าเว็บของสำนักข่าวต้นทางอาจมีการปรับเปลี่ยน (DOM Mutation) ซึ่งส่งผลให้โมดูลแยกวิเคราะห์ (Parser) ต้องได้รับการปรับปรุงให้สอดคล้องกับโครงสร้างหน้าเว็บใหม่อย่างต่อเนื่อง นอกจากนี้ การสรุปเนื้อหาข่าวยังพึ่งพาการเชื่อมต่อกับบริการ API ภายนอก ได้แก่ KKU Gen.AI หรือ Large Language Model API ซึ่งในกรณีที่เกิดปัญหาการสื่อสารบนเครือข่ายหรือโควตาการเรียกใช้งานหมดลง ระบบจะปรับเปลี่ยนการทำงานเข้าสู่โหมดการแสดงผลข้อความสรุปตั้งต้นแทนโดยอัตโนมัติ

\subsection{ข้อจำกัดด้านทรัพยากรและการประมวลผล}

ในมิติด้านทรัพยากรระบบ การเปิดใช้งาน Playwright Headless Browser เพื่อสกัดข้อมูลจากเว็บไซต์ข่าวสารที่มีระบบป้องกันบอตหรือใช้การแสดงผลฝั่งไคลเอนต์ จำเป็นต้องใช้ทรัพยากรหน่วยความจำหลักและรอบการประมวลผลของหน่วยประมวลผลกลาง (CPU) สูงกว่าการดึงข้อมูลผ่านโพรโทคอล HTTP REST API ปกติอย่างมีนัยสำคัญ

% -----------------------------------------------------------
\section{ปัญหา อุปสรรค และแนวทางแก้ไข}
% -----------------------------------------------------------

ปัญหาที่พบในระหว่างการพัฒนาและแนวทางที่ใช้แก้ไขแสดงดังตารางที่ \ref{tab:problems}

\begin{table}[H]
\centering
\caption{ปัญหา อุปสรรค และแนวทางการแก้ไข}
\label{tab:problems}
\resizebox{\textwidth}{!}{%
\begin{tabular}{|c|p{5.5cm}|p{6.5cm}|}
\hline
\textbf{ลำดับ} & \textbf{ปัญหาที่พบ} & \textbf{แนวทางแก้ไข} \\
\hline
1 & บางเว็บไซต์ข่าวมีระบบป้องกันบอต (Cloudflare/Bot detection) ทำให้ไม่สามารถดึงข้อมูลผ่าน HTTP Client ธรรมดาได้ & ใช้ Playwright Browser ควบคุมเบราว์เซอร์จริงในการเรนเดอร์หน้าเว็บและจัดการ Cookie Session \\ \hline
2 & การเรียก LLM API เพื่อสรุปข่าวทุกครั้งทำให้การโหลดข้อมูลล่าช้าและสิ้นเปลืองโควตา & ออกแบบระบบ Multi-tier Cache (\texttt{AsyncInMemoryCache} และ Local JSON Persistence) เพื่อเก็บผลลัพธ์การสรุป \\ \hline
3 & ข่าวสารภาษาไทยไม่มีการเว้นวรรคระหว่างคำ ทำให้การจับคู่หมวดหมู่มีความคลาดเคลื่อน & ใช้ PyThaiNLP ในการตัดคำแบบพจนานุกรมและคัดกรอง Stop words ก่อนนำไปจับคู่หมวดหมู่ \\ \hline
4 & ข่าวเรื่องเดียวกันจากต่างสำนักพิมพ์มีมุมมองการพาดหัวต่างกัน และมีการกำหนด Tag หมวดหมู่ขัดแย้งกัน & ใช้การวัดความคล้ายคลึงแบบไฮบริด (หัวข้อ, สรุปย่อ, เอนทิตีบุคคล/สถานที่, คำสำคัญ) ร่วมกับกลไกมติเสียงข้างมาก (Majority Voting) ในระดับกลุ่มข่าว \\
\hline
\end{tabular}%
}
\end{table}

% -----------------------------------------------------------
\section{ข้อเสนอแนะในการพัฒนาต่อไป}
% -----------------------------------------------------------

\subsection{การปรับปรุงฟังก์ชันการทำงาน}

สำหรับการพัฒนาฟังก์ชันการทำงานในระยะต่อไป ระบบสามารถขยายขีดความสามารถได้หลายด้าน ประการแรกคือการพัฒนา Adapter เชื่อมต่อกับระบบจัดการฐานข้อมูลแบบกระจายศูนย์ (Distributed Database) เช่น MongoDB หรือ PostgreSQL เพิ่มเติมจาก \texttt{FileNewsRepository} เพื่อรองรับการจัดเก็บและสืบค้นคลังข่าวสารขนาดใหญ่ได้อย่างมีเสถียรภาพ ประการถัดมาคือการพัฒนาระบบสมาชิกและบัญชีผู้ใช้งานเพื่ออำนวยความสะดวกในการซิงค์ข้อมูลข่าวที่คั่นหน้า (Bookmarks) และประวัติความสนใจข้ามอุปกรณ์ และประการสุดท้ายคือการพัฒนาระบบปรับค่าน้ำหนัก $w_1, \dots, w_4$ ของอัลกอริทึมจัดกลุ่มข่าวแบบอัตโนมัติตามพฤติกรรมการอ่านจริงของผู้ใช้งาน

\subsection{แนวทางการวิจัยต่อยอด}

ในเชิงการวิจัยต่อยอด คณะผู้พัฒนาเล็งเห็นแนวทางสำคัญ 3 ประการ ประการแรกคือการประยุกต์ใช้โมเดล Transformer ภาษาไทยเฉพาะทางและ Cross-Encoder เช่น WangchanBERTa มาคำนวณคะแนนความคล้ายคลึงเชิงความหมาย (Semantic Similarity) เพื่อยกระดับความแม่นยำในการจัดกลุ่มข่าวสารที่มีสำนวนภาษาต่างกันอย่างลึกซึ้ง ประการที่สองคือการดำเนินการทดสอบประเมินผลกับชุดข้อมูลข่าวจริงภาคสนาม (Real-World Clustering Benchmark) โดยสร้างชุดข้อมูลข่าวข้ามสำนักข่าวพร้อมป้ายกำกับเหตุการณ์ เพื่อวัดค่า Pairwise Precision, Recall และ F1-score ตามแผนงานในอนาคต และประการที่สามคือการพัฒนาโมดูลตรวจจับและเปรียบเทียบมุมมองข่าวสารข้ามสำนักข่าว (Cross-source Fact Verification) เพื่อวิเคราะห์ความถูกต้องและความเป็นกลางของเนื้อหาข่าวสารในสังคมต่อไป

